\subsection{Exact Arc-Intersection Algorithm (Benchmark Implementation)}
\label{app:arc-algorithm}

This subsection provides a concrete, implementation-ready algorithm for computing the overlap lengths
$\tilde L_{PP},\tilde L_{PM},\tilde L_{MP},\tilde L_{MM}$ on the circle $\phi\in[0,\pi)$.
It is designed to be numerically robust (handles wrap-around cleanly) and directly supports benchmarking
against simulation.

\paragraph{Arc representation.}
Represent any arc (interval on the circle) by its endpoints $(s,e)$ with $s,e\in[0,\pi)$ and
a convention that the arc includes points encountered when traveling \emph{forward} from $s$ to $e$
modulo $\pi$.
If $s\le e$, the arc is the ordinary interval $[s,e]$; if $s>e$, it wraps and denotes $[s,\pi)\cup[0,e]$.

\paragraph{Linearization of wrapped arcs.}
For intersection computations, convert each wrapped arc into a union of at most two non-wrapping intervals on $[0,\pi)$:
\[
\mathrm{split}(s,e)=
\begin{cases}
\{[s,e]\}, & s\le e,\\
\{[s,\pi),\ [0,e]\}, & s>e.
\end{cases}
\]

\paragraph{Arc length.}
For a non-wrapping interval $[s,e]$ with $0\le s\le e\le \pi$, length is $e-s$.
For a wrapped arc $(s,e)$, length is $(\pi-s)+e$.

\paragraph{Centered arcs.}
For a center $\mu\in[0,\pi)$ and half-width $m\in[0,\pi/2]$, define the wrapped arc
\[
\mathrm{arc}(\mu,m) = (\mu-m,\ \mu+m)\ \ \ (\text{mod }\pi),
\]
i.e.\ endpoints reduced modulo $\pi$.

Threshold acceptance arcs are
\[
P_\theta=\mathrm{arc}(\theta,m),\qquad
M_\theta=\mathrm{arc}\!\left(\theta+\frac{\pi}{2},m\right),
\]
with $m=m_{\mathrm{thr}}(t)$.

\paragraph{Coincidence arcs.}
Assume the coincidence acceptance set $C_{\theta_A,\theta_B}(\Delta)$ has been decomposed into a finite union
\[
C_{\theta_A,\theta_B}(\Delta)=\bigcup_{k=1}^{K} I_k,
\]
where each $I_k$ is a wrapped arc. (This decomposition can be obtained analytically in symmetric cases,
or numerically by solving for boundaries in $\phi$ where
$|t_c(u_A(\phi))-t_c(u_B(\phi))|=\Delta$.)

Then the modified overlap length for a pair of acceptance arcs $X\in\{P_{\theta_A},M_{\theta_A}\}$ and
$Y\in\{P_{\theta_B},M_{\theta_B}\}$ is
\[
\tilde L_{XY}=\sum_{k=1}^{K} |X\cap Y\cap I_k|.
\]

%------------------------------------------------------------
\subsubsection*{Algorithmic Pseudocode}

Below we give pseudocode using \texttt{algorithmic} (or \texttt{algorithmicx}) style. The functions are:
\begin{itemize}
    \item \textsc{SplitArc}$(s,e)$: returns list of non-wrapping intervals on $[0,\pi)$.
    \item \textsc{IntersectIntervals}$(I,J)$: returns intersection interval (possibly empty) of non-wrapping intervals.
    \item \textsc{Length}$(I)$: returns length of non-wrapping interval.
    \item \textsc{Arc}$(\mu,m)$: returns wrapped arc endpoints $(s,e)$ reduced modulo $\pi$.
    \item \textsc{Overlap3}$(A,B,C)$: computes length of intersection of three wrapped arcs.
\end{itemize}

\begin{algorithmic}[1]
\Function{ModPi}{$x$}
    \State \Return $x - \pi\lfloor x/\pi \rfloor$ \Comment{maps $x$ into $[0,\pi)$}
\EndFunction

\Function{Arc}{$\mu, m$}
    \State $s \gets$ \Call{ModPi}{$\mu - m$}
    \State $e \gets$ \Call{ModPi}{$\mu + m$}
    \State \Return $(s,e)$ \Comment{wrapped arc from $s$ to $e$ mod $\pi$}
\EndFunction

\Function{SplitArc}{$(s,e)$}
    \If{$s \le e$}
        \State \Return $\{[s,e]\}$
    \Else
        \State \Return $\{[s,\pi], [0,e]\}$
    \EndIf
\EndFunction

\Function{IntersectIntervals}{$[a,b],[c,d]$}
    \State $p \gets \max(a,c)$
    \State $q \gets \min(b,d)$
    \If{$p \ge q$}
        \State \Return $\emptyset$
    \Else
        \State \Return $[p,q]$
    \EndIf
\EndFunction

\Function{Length}{$[a,b]$}
    \State \Return $b-a$
\EndFunction

\Function{Overlap3}{$A,B,C$}
    \Comment{$A,B,C$ are wrapped arcs given by endpoints $(s,e)$}
    \State $\mathcal{A}\gets$ \Call{SplitArc}{$A$}
    \State $\mathcal{B}\gets$ \Call{SplitArc}{$B$}
    \State $\mathcal{C}\gets$ \Call{SplitArc}{$C$}
    \State $L\gets 0$
    \ForAll{$I\in\mathcal{A}$}
        \ForAll{$J\in\mathcal{B}$}
            \State $IJ\gets$ \Call{IntersectIntervals}{$I,J$}
            \If{$IJ \neq \emptyset$}
                \ForAll{$K\in\mathcal{C}$}
                    \State $IJK\gets$ \Call{IntersectIntervals}{$IJ,K$}
                    \If{$IJK \neq \emptyset$}
                        \State $L\gets L +$ \Call{Length}{$IJK$}
                    \EndIf
                \EndFor
            \EndIf
        \EndFor
    \EndFor
    \State \Return $L$
\EndFunction

\Function{ComputeOverlaps}{$m,\theta_A,\theta_B,\{I_k\}_{k=1}^K$}
    \State $P_A\gets$ \Call{Arc}{$\theta_A, m$}
    \State $M_A\gets$ \Call{Arc}{$\theta_A+\pi/2, m$}
    \State $P_B\gets$ \Call{Arc}{$\theta_B, m$}
    \State $M_B\gets$ \Call{Arc}{$\theta_B+\pi/2, m$}
    \State $\tilde L_{PP}\gets 0,\ \tilde L_{PM}\gets 0,\ \tilde L_{MP}\gets 0,\ \tilde L_{MM}\gets 0$
    \For{$k=1$ to $K$}
        \State $C\gets I_k$
        \State $\tilde L_{PP}\gets \tilde L_{PP} +$ \Call{Overlap3}{$P_A,P_B,C$}
        \State $\tilde L_{PM}\gets \tilde L_{PM} +$ \Call{Overlap3}{$P_A,M_B,C$}
        \State $\tilde L_{MP}\gets \tilde L_{MP} +$ \Call{Overlap3}{$M_A,P_B,C$}
        \State $\tilde L_{MM}\gets \tilde L_{MM} +$ \Call{Overlap3}{$M_A,M_B,C$}
    \EndFor
    \State \Return $(\tilde L_{PP},\tilde L_{PM},\tilde L_{MP},\tilde L_{MM})$
\EndFunction
\end{algorithmic}

\paragraph{From overlaps to correlations.}
Given $(\tilde L_{PP},\tilde L_{PM},\tilde L_{MP},\tilde L_{MM})$, compute
\begin{align}
\tilde D &= \tilde L_{PP}+\tilde L_{PM}+\tilde L_{MP}+\tilde L_{MM},\\
\tilde N &= (\tilde L_{PP}+\tilde L_{MM})-(\tilde L_{PM}+\tilde L_{MP}),\\
\mathcal{E}^{(\mathrm{coin})} &= \tilde N/\tilde D.
\end{align}
Then compute CHSH $S$ from the four setting pairs.

\paragraph{Notes for benchmarking.}
\begin{itemize}
    \item The circle circumference is $\pi$ because $\phi$ and $\phi+\pi$ define the same linear polarization.
    \item The algorithm is exact given a correct decomposition $\{I_k\}$ of the coincidence set.
    \item If coincidence selection is disabled, set $K=1$ and $I_1=(0,\pi)$, recovering the threshold-only correlations.
\end{itemize}
